%% ma: L11-B14-0004
%% lop: 11
%% chuong: 4
%% bai: 14
%% dang: 11.14.4
%% loai: TL
%% muc: VD
%% dap_an: "-"
%% nguon: "tai-lieu-hoc-tap-mon-toan-lop-11-theo-chuong-trinh-giao-duc-pho-thong-2018.pdf (Lớp Toán Thầy Duy, Phạm Lê Duy, THPT Chu Văn An), Chương 4 Bài 5, Dạng 3, Ví dụ 4 (tr.524)"
%% kien_thuc: "Lớp 11 Bài 14 (phép chiếu song song); tọa độ/vectơ trong không gian được dùng để kiểm định (chương 7 chưa học); cách làm gốc dựa vào phép chiếu song song theo phương BN hoặc tương đương"
%% trang_thai: chinh-thuc
%% ghi_chu: "Đáp số 1 do kho tự giải (tài liệu gốc không có lời giải mẫu) và đã kiểm bằng tọa độ sympy. Lời giải dùng tọa độ (kiến thức Toán 12); lời giải thuần túy dùng phép chiếu song song cần giáo viên bổ sung. Đã duyệt 03/11/2026."
\begin{ex}
Cho hình hộp $ABCD.A^{\prime}B^{\prime}C^{\prime}D^{\prime}$. Gọi $M,N$ lần lượt là trung điểm của $CD$ và $CC^{\prime}$. Gọi $\Delta$ là đường thẳng đi qua $M$, đồng thời cắt $AN$ và $A^{\prime}B$. Gọi $I,J$ lần lượt là giao điểm của $\Delta$ với $AN$ và $A^{\prime}B$. Hãy tính tỉ số $\dfrac{IM}{IJ}$.
\loigiai{\sloppy Chọn hệ trục với $A$ là gốc, $\overrightarrow{AB},\overrightarrow{AD},\overrightarrow{AA^{\prime}}$ là các vectơ đơn vị (phép chiếu song song và tỉ số độ dài không phụ thuộc hình dạng hộp). Khi đó $M\left(\dfrac12;1;0\right)$, $N\left(1;1;\dfrac12\right)$, $A^{\prime}(0;0;1)$, $B(1;0;0)$. Điểm $I\in AN$: $I=(t;t;\tfrac t2)$; điểm $J\in A^{\prime}B$: $J=(s;0;1-s)$. Ba điểm $M,I,J$ thẳng hàng cho $t=s=\dfrac12$, nên $I=\left(\dfrac12;\dfrac12;\dfrac14\right)$ và $J=\left(\dfrac12;0;\dfrac12\right)$. Khi đó $I=\dfrac{M+J}2$, tức $I$ là trung điểm của $MJ$, nên $\dfrac{IM}{IJ}=1$.}
\end{ex}
